\section*{Capítulo 2: Los padres de la gravedad}
\begin{description}[style=multiline,leftmargin=2em,labelwidth=1.5em]

% DOCX paragraph 2017
\item[1.]\hypertarget{note:02-01}{}En un sistema de referencia en el que los dos
cuerpos tienen una posición dada por los vectores\index{vector}
\({\overrightarrow{r}}_{1}\) y \({\overrightarrow{r}}_{2}\), su separación viene
dada por la diferencia vectorial entre ambos,
\({\overrightarrow{r} = \overrightarrow{r}}_{1} - {\overrightarrow{r}}_{2}\),
cuya magnitud, o módulo, se indica como
\(\left| \overrightarrow{r} \right| =\) \emph{r}.

% DOCX paragraph 2018
\item[2.]\hypertarget{note:02-02}{}El segundo de arco es una unidad de medida de
ángulos que se usa tanto en navegación como en astronomía y representa una manera
alternativa de expresar la amplitud del ángulo abarcado por un arco. En la práctica,
además de medir la amplitud de un ángulo en grados, unas unidades con las que
estamos familiarizados pero que no resultan prácticas en navegación ni en
astronomía, se introducen unidades más pequeñas: los \emph{minutos} de arco y
los \emph{segundos} de arco. Así pues, un arco de un grado se divide en
60 minutos, cada uno de los cuales corresponde a 0,017 grados, o en 3600 segundos,
cada uno de los cuales corresponde a 0,00028 grados. Si usáramos grados en vez
de arcosegundos, ¡diríamos que el perihelio\index{perihelio} de Mercurio varía
tan solo 0,0121 grados por siglo!

\end{description}
